\section{Introduction}

The revolutionary success of Large Language Models (LLMs) at generating human-like texts~\citep{brown2020language,bubeck2023sparks,chowdhery2023palm} has raised several societal concerns regarding the misuse of LLM outputs. Unregulated LLM outputs pose risks ranging from the contamination of future training corpora~\citep{shumailov2023curse,das2024under} to the propagation of disinformation~\citep{zellers2019defending,vincent2022ai} and academic misconduct~\citep{jarrah2023using,milano2023large}. Consequently, there is an urgent need for reliable detection mechanisms capable of distinguishing LLM-generated text from human-authored content.

To address this challenge, \citet{aaronson2022my} and \citet{kirchenbauer2023watermark} propose \emph{statistical watermarking} as a theoretically-grounded method to distinguish machine-generated content from a target distribution. These methods operate by injecting statistical signals into the generated texts in the decoding phase of LLMs. Mechanistically, the watermark generator couples the output tokens by a sequence of pseudorandom seeds. During detection, the detector reconstructs the seeds and performs a hypothesis test to check the dependence between the seeds and the observed tokens. This formulation allows the detection problem to be treated as a rigorous statistical hypothesis test~\citep{huang2023towards} where the detection power and sample complexity are governed by the randomness in the target distribution. 
Various designs for the underlying seed distributions have been explored, such as binary~\citep{christ2023undetectable}, exponential~\citep{kuditipudi2023robust}, and Gaussian~\citep{block2025gaussmark}.  Recently,  \citet{huang2025watermarking} advanced the concept of using an auxiliary open-source model as an \emph{anchor} to generate random
seeds and coupling them with tokens via speculative decoding ~\citep{leviathan2023fast}. This approach is motivated by the theoretical insight that watermarking achieves the optimal power when the seed distribution approximates the target distribution~\citep{huang2023towards}. By exploiting the distributional proximity between modern LLMs, anchor-based methods have shown promise in improving both detection efficiency
and robustness to removing-attacks~\citep{piet2023mark}. 
Empirically, watermarked texts can be detected within short token horizons ($<100$ tokens) and high success rate (nearly $50\%$ under paraphrasing) under various attacks.

Despite these successes, state-of-the-art watermarking remains constrained by two fundamental limitations. First, the design of generation and detection schemes lacks a unifying guiding principle.  \citet{huang2025watermarking} utilize a hashed green-red-list seed adapted from \citet{kirchenbauer2023watermark}, but this heuristic lacks rigorous justification. What is needed is a characterization of optimal generation and detection schemes under given anchor distributions.
Second, there exists a methodological disconnect between generation and detection. While text generation is inherently autoregressive and variable in length, current detection paradigms are \emph{fixed-horizon} and \emph{ad hoc}. In practice, for example, a detector may wish to process a stream of tokens and stop as soon as confidence is high. However, with the standard statistical paradigm based on p-values, it is not allowed to continuously monitor the test statistic and decide when to stop based on the current value---this is known as ``p-hacking'' and it invalidates the Type-I error guarantee. This inability to perform \emph{post-hoc} sequential analysis severely limits detection efficiency. These gaps highlight a critical research challenge:
\begin{center}
    \emph{How can we design provably efficient watermarking schemes leveraging anchor distributions that allow for valid early stopping?}
\end{center}

In this work, we provide the first formal resolution to this question. We formulate the problem of \emph{Anchored E-Watermarking}, where both the generator and detector share access to an anchor distribution $p_0$, aiming to watermark a family of distributions in the $\delta$-proximity of $p_0$. Departing from classical p-value analysis, we adopt \emph{e-values} as the central detection paradigm. E-values~\citep{vovk1993logic,shafer2011testmartingales,vovk2021values} are nonnegative random variables $E$ satisfying $\E_{\mathbf{H}_0}[E] \leq 1$ under the null. Unlike p-values, e-values arise from supermartingales, and it is therefore possible to preserve Type-I error guarantees under optional stopping based on ongoing analysis of data. We analyze the optimal e-value for watermark detection with respect to the worst-case log-growth rate~\citep{kelly1956new} and the expected stopping time~\citep{grunwald2020safe,waudby2025universal}, thus fully characterize the average per-step growth rate of evidence and sample efficiency. Our main theoretical contributions are summarized informally as follows:

\begin{theorem}[Informal version of \Cref{thm:log-growth} and \Cref{thm:stopping-time}]\label{thm:informal}
The optimal worst-case log-growth rate $\E_{\mathbf{H}_1}[\log E]$ under the alternative $\mathbf{H}_1$ is given by:
\begin{align*}
    J^* = h + \left(1-\frac{\delta}{2}\right)\log\left(1-\frac{\delta}{2}\right)
 + \frac{\delta}{2}\log\left(\frac{\delta}{2(n-1)}\right),
\end{align*}
where $h = H(p_0)$ is the Shannon entropy of the anchor distribution $p_0$ and $\delta>0$ is a robustness tolerance parameter.
Furthermore, the optimal expected stopping time to achieve a Type-I error $\alpha$ scales as $\frac{\log(1/\alpha)}{J^*}$.
\end{theorem}

To the best of our knowledge, this work represents the first application of e-values to the domain of statistical watermarking. By enabling valid sequential testing, our framework significantly improves detection efficiency, allowing the system to flag machine-generated text with fewer tokens than fixed-horizon counterparts. With the ability to early-stop, this sequential paradigm enhances robustness against adaptive attacks: because the detector can terminate immediately upon accumulating sufficient evidence, the watermark remains effective even if the attacker perturbs the text heavily in later segments (post-stopping). Consequently, our approach offers a theoretically rigorous and practically superior alternative to existing heuristic detection methods. Through experiments on real watermarking benchmark~\citep{piet2023mark}, we show that the theoretical scheme implied by \Cref{thm:informal} achieves consistently higher sample efficiency than state-of-the-art methods, reducing the token consumption by 13-15\% across various temperature settings.

\subsection{Related works}

\label{sec:related_work}

\paragraph{Statistical watermarking.}
Watermarking offers a white-box provenance mechanism for detecting LLM-generated text~\citep{tang2023science}, complementing post-hoc detectors and provenance tools developed for neural text generation~\citep{zellers2019defending}.
Classical digital watermarking and steganography provide a broad toolbox for embedding and extracting imperceptible signals under benign or adversarial channel edits~\citep{cox2007digital}. Early works in NLP literature studied watermarking and tracing of text via editing or synonym substitutions~\citep{venugopal-etal-2011-watermarking,rizzo2019fine,abdelnabi2021adversarial,yang2022tracing,kamaruddin2018review}.
In contrast, modern \emph{statistical} (a.k.a.\ generative) watermarking~\citep{aaronson2022my,kirchenbauer2023watermark} injects a secret, testable distributional bias into the sampling process, and detects this bias via hypothesis testing on the generated token sequence.
A rapidly growing theory studies efficiency and optimality of watermarking tests and encoders: results include finite-sample guarantees, information-theoretic limits, and constructions that are distortion-free or unbiased~\citep{huang2023towards,zhao2023provable,li2024statistical,block2025gaussmark,kuditipudi2023robust,hu2023unbiased,xie2025debiasing}.
Negative and hardness results highlight fundamental limitations against adaptive or distribution-matching adversaries~\citep{christ2023undetectable,christ2024pseudorandom,golowich2024edit}, motivating alternative design considerations such as distribution-preserving and public-key  schemes~\citep{wu2023distributionpreserving,liu2023unforgeable,fairoze2023publiclydetectable}.
Empirical robustness is commonly assessed under paraphrasing, editing and translation, with recent work studying cross-lingual failure modes and defenses~\citep{he2024translation}, and benchmarks/frameworks such as MarkMyWords and scalable pipelines for watermark evaluation and deployment~\citep{piet2023mark,zhang2024remarkllm,lau2024waterfall,dathathri2024scalable}.
Finally, a complementary line of work leverages semantic structure and auxiliary models to boost detection power under benign distributional structure, including semantic/paraphrastic watermarks and speculative-sampling-based schemes~\citep{ren2024robustsemanticsbasedwatermarklarge,liu2024adaptive,hou2024k,hou2024semstampsemanticwatermarkparaphrastic,fu2024watermarking,huang2025watermarking,leviathan2023fast}.
Our anchored e-value approach builds on these works by explicitly treating detection as anytime-valid sequential testing~\citep{chen2024online}, and by allowing model-assisted calibration under distribution shift~\citep{huang2025watermarking,he2024distributionadaptive,he2025distributionalembedding}.

\paragraph{E-values and sequential hypothesis testing.}
Sequential hypothesis testing studies inference procedures that remain valid under data-dependent stopping and streaming data~\citep{WaldSequentialAnalysis1947,Robbins1952SomeAO,breimualphanu2011optimal}.
Classical likelihood-ratio and $p$-value based methods can fail under optional stopping or composite hypotheses, motivating the use of nonnegative supermartingales / test martingales whose stopped values are valid evidence measures~\citep{shafer2011testmartingales,vovk2021values,Shafer2021Testing,grunwald2020safe,ramdas2023game}.
The resulting \emph{e-values} and \emph{e-processes} unify betting scores, likelihood ratios, and (stopped) Bayes factors. Additionally, e-values and e-processes support modular operations such as merging and calibration~\citep{vovk2021values,ramdasWang2025hypevalues}.
This framework has enabled principled notions of power and optimality, including log-/growth-rate criteria and sharp asymptotic rates for broad classes of betting-based tests~\citep{grunwald2020safe,waudby2025universal}.
E-values also interact fruitfully with multiple testing and online testing, yielding e-value analogues of the BH and wealth-based procedures and their extensions to structured settings~\citep{wangramdas2022fdrevalues,ramdas2017decayingmemory,ramdas2018saffron,ramdas2019dagger}.
Recent work further broadens the reach of e-values to new ML settings, such as conformal prediction and sequential monitoring of strategic systems~\citep{gauthier2025conformal_evalues,gauthier2026bettingequilibrium,aolaritei2025sgdcs}, and connects e-values to model-assisted efficiency gains via prediction-powered inference~\citep{wasserman2020universal,csillag2025predictionpoweredevalues}.
Our work draws on these foundations to construct e-values tailored to watermark detection which ensures validity under arbitrary stopping rules and enables sequential evidence accumulation in practical provenance audits.

